\documentclass[11pt]{article}
\usepackage[utf8]{inputenc}
\usepackage{amsmath,amssymb}
\usepackage[margin=1in]{geometry}
\usepackage{hyperref}
\emergencystretch=3em

\title{The chart integral as a genuine $d$-dimensional integral:\\
$\int_{(0,b]^d}\prod u_i^{h_i}\,f\big(c\prod u_i^{k_i}\big)\,du=\texttt{iterChartGen}\,d\,c$}
\author{Claude, with Daniel Murfet}
\date{2026-09-06}

\begin{document}
\maketitle
\sloppy

\begin{abstract}
The recursive \texttt{iterChartGen} of unit~50 integrates one coordinate at a time; the paper writes
the standard integral over the box $(0,b]^d\subseteq\mathbb R^d$. This unit defines the box integral as
a Bochner integral on $\mathrm{Fin}\,d\to\mathbb R$ against the product of the restricted Lebesgue
measures and proves it equal to the recursion, so that the multiplicity theorems of units~50 and~52
hold verbatim in paper form (\texttt{boxIntegral\_isEquivalent},
\texttt{boxIntegral\_mixed\_isEquivalent}). Zero \texttt{sorry}/\texttt{axiom}.
\end{abstract}

\section{Setting}
Splitting off the last coordinate of $\mathrm{Fin}(d+1)\to\mathbb R$ is the measurable equivalence
\texttt{piFinSuccAbove} at \texttt{Fin.last d}, which is measure preserving from the product of the
$d+1$ restricted measures to the product of one restricted measure with the product of the other $d$
(\texttt{measurePreserving\_piFinSuccAbove}). Fubini on this product reduces the $(d+1)$-fold integral
to $\int_0^bu^{h_d}\,(\text{$d$-fold integral at scale }cu^{k_d})\,du$, which is the recursion.

\section{The things formalised}
{\small
\begin{verbatim}
noncomputable def boxMeasure (b : ℝ) (d : ℕ) : Measure (Fin d → ℝ) :=
  Measure.pi fun _ => (volume : Measure ℝ).restrict (Ioc 0 b)
noncomputable def boxIntegral (β a b c : ℝ) (k h : ℕ → ℕ) (d : ℕ) : ℝ :=
  ∫ u : Fin d → ℝ, (∏ i, u i ^ h i) * quadKernel β a (c * ∏ i, u i ^ k i) ∂(boxMeasure b d)
theorem boxIntegral_eq_iterChartGen (β a b : ℝ) (k h : ℕ → ℕ) (d : ℕ) (c : ℝ) :
    boxIntegral β a b c k h d = iterChartGen β a b k h d c
theorem boxIntegral_isEquivalent ... :
    (fun n => boxIntegral β a b (Real.sqrt n) k h (d + 2)) ~[atTop]
      fun n => (∏ i ∈ Finset.range (d + 2), (1 : ℝ) / k i)
        * (weightedMass β a (p - 1) / (2 ^ (d + 1) * ((d + 1).factorial : ℝ)))
        * (n ^ (-(p / 2)) * Real.log n ^ (d + 1))
theorem boxIntegral_mixed_isEquivalent ...
\end{verbatim}
}

\section{Proof strategy}
Induction on $d$ generalising the scale $c$. For $d=0$ the product measure over the empty index is a
Dirac measure (\texttt{Measure.pi\_of\_empty}) and the integral is the kernel at $c$. For the step,
the integrand is rewritten as $g\circ e$ with $g(t,v)=(\prod_jv_j^{h_j})t^{h_d}\,f(c\prod_jv_j^{k_j}t^{k_d})$
using \texttt{Fin.prod\_univ\_castSucc} and the two definitional facts $(eu).1=u(\mathrm{last})$,
$(eu).2\,j=u(\mathrm{castSucc}\,j)$ (\texttt{Fin.succAbove\_last\_apply}); then
\texttt{MeasurePreserving.integral\_comp'} and \texttt{integral\_prod}. Integrability of $g$ on the
product of restricted measures is obtained by identifying that measure with Lebesgue measure restricted
to the box (\texttt{Measure.prod\_restrict}, \texttt{Measure.restrict\_pi\_pi}) and using continuity on
the compact $[0,b]\times[0,b]^d$. The inner integral is the induction hypothesis at scale $cu^{k_d}$
after pulling out $u^{h_d}$.

\section{Roadblocks and resolutions}
\begin{itemize}
\item \texttt{simp} would not unfold the structure-literal coercion of \texttt{piFinSuccAbove}; the
two component identities are \texttt{rfl} up to \texttt{succAbove\_last\_apply}, so state them as
\texttt{have}s and \texttt{simp only} with them under the product binder.
\item \texttt{fun\_prop} does not know \texttt{quadKernel}; assemble continuity by hand from
\texttt{continuous\_finsetProd}, \texttt{continuous\_apply}, and \texttt{quadKernel\_continuous}.
\item \texttt{Measure.volume\_pi} is not the name on this pin; since \texttt{volume} on a pi type is
definitionally \texttt{Measure.pi}, a \texttt{change} to the \texttt{Measure.pi} form suffices before
\texttt{restrict\_pi\_pi}.
\end{itemize}

\section{What was Mathlib, what was new}
Mathlib: \texttt{measurePreserving\_piFinSuccAbove}, \texttt{MeasurePreserving.integral\_comp'},
\texttt{integral\_prod}, \texttt{Measure.pi\_of\_empty}, \texttt{integral\_dirac},
\texttt{Measure.restrict\_pi\_pi}, \texttt{Fin.prod\_univ\_castSucc}, \texttt{isCompact\_univ\_pi}. New:
the box integral and its identification with the recursion.

\section{Lessons}
Keeping the recursion as the working definition and proving the paper-form identity once at the end
was the right order: all analysis was done on one-dimensional objects, and the $d$-dimensional
measure theory is confined to a single induction.

\section{Outcome}
The multiplicity theorems for the chart standard integral with arbitrary $(k,h)$ are now stated on
$(0,b]^d\subseteq\mathbb R^d$ exactly as in the paper, in the all-equal and one-noncritical regimes.
\end{document}
